% GENERATED FILE. Edit canonical lessons, then run npm run generate:books.
% canonical-source-hash: fnv1a64:0438ff201af1c6f0
% canonical-lessons: ZH-C62-kuai, ZH-W62-kuai-fast, ZH-C62-jiu, ZH-C62-yuan, ZH-C62-jin, ZH-C62-zhidao

\chapter{Fast, Far and Near}
\label{ch:zh-fast-far-and-near}
\hlchaptermodality{62}

\begin{chapteropening}
\textbf{By the end of this chapter:} \emph{I can say fast, old, far and near, and say that I know something.}

Complete the last lesson of chapter 62: I can say fast, old, far and near, and say that I know something.
\end{chapteropening}

\section[\texorpdfstring{\zh{快} (kuài)}{快 (kuài)}]{\zh{快} (kuài) — fast; quick}

\label{lesson:ZH-C62-kuai}



\begin{quote}
Before the new one: say the Mandarin for short in height, then the Mandarin for new.
\end{quote}

\subsection*{You'll want to know: \zh{快}}
\textbf{\zh{快}} — \emph{kuài} — \textquotedblleft{}fast; quick\textquotedblright{}.

\begin{quote}
\textbf{\zh{我的手机很快}} \emph{wǒ de shǒujī hěn kuài} — My phone is fast.
\end{quote}

The new character, in pieces: \zh{快} — on the left the heart piece from the foot of \zh{想}, standing upright, and on the right a four-stroke piece read guài. The next lesson puts it in your hand.

\begin{cousinweb}
\textbf{\zh{快}} (\emph{kuài}) is inherited from Middle Chinese \emph{khwaejH} (Guangyun, in Baxter's transcription), glossed 'to one's liking' and 'glad'. The sense 'fast' is not given there, so how it got it is not settled.
\end{cousinweb}

\begin{grammarlens}[title={What you've built}]
One more word for everyday answers.
\end{grammarlens}

\subsection*{Your turn}
\begin{itemize}
\raggedright
  \item \emph{Say it:} \emph{kuài}
  \item \emph{Say it:} \emph{kuài}, once more
  \item \emph{Say it:} \emph{wǒ de shǒujī hěn kuài} — My phone is fast.
  \item \emph{Recall:} say the Mandarin for short in height, then the Mandarin for new, then say \emph{kuài} again
\end{itemize}

\begin{tcolorbox}[breakable,colback=teal!4,colframe=teal!35!black,title={Before you move on}]
What does \zh{快} mean? (\textquotedblleft{}Fast; quick\textquotedblright{}.) Say it once more.
\end{tcolorbox}
\section[\texorpdfstring{\zh{快} (kuài)}{快 (kuài)}]{\zh{快} — the upright heart, then guài}

\label{lesson:ZH-W62-kuai-fast}



\begin{quote}
The word first, then the character inside it.

\begin{itemize}
\raggedright
  \item \emph{Recall:} say \emph{kuài} — \textbf{R1}, one lesson back, before it decays
\end{itemize}

\begin{quote}
\zh{想} \zh{大}
\end{quote}

\emph{Read it:} each one aloud, with nothing said first — \emph{xiǎng}, \emph{dà}
\end{quote}

\subsection*{Script — the upright heart, then guài}
\hlblockfigure{\detokenize{figures/ZH-W62-kuai-fast-filmstrip.pdf}}{How \zh{快} is written, stroke by stroke}

\begin{quote}
\zh{快}
\end{quote}

\textbf{Two pieces, left before right.}

\textbf{The upright heart piece, three strokes}: a short dot on the left, a short dot on the right, then a long vertical down between them: the heart piece from the foot of \zh{想}, standing upright.

\textbf{The right piece, four strokes, read guài}: a horizontal that turns down at its right end, a long horizontal across, a long left-falling stroke from the top down through the bars, and a stroke falling to the right.

3 + 4 = seven.

\subsection*{Your turn}
\begin{itemize}
\raggedright
  \item \emph{Write it:} \zh{快}, all seven, with the model in front of you
  \item \emph{Cover:} write it again from memory, and count the pieces, not the strokes
\end{itemize}

\begin{tcolorbox}[breakable,colback=teal!4,colframe=teal!35!black,title={Before you move on}]
How many pieces in \zh{快}? (\textbf{Two}.) Which dot comes first? (\textbf{The left one}.) How many strokes on the right side? (\textbf{Four}.)

Source: \href{https://raw.githubusercontent.com/chanind/hanzi-writer-data/68d10a4b21150cae5e1ebbd223eed289cf32d90c/data/\%E5\%BF\%AB.json}{Hanzi Writer Data \zh{快}, pinned stroke paths and medians}.
\end{tcolorbox}
\section[jiù]{jiù — old (of things); no longer new}

\label{lesson:ZH-C62-jiu}



\begin{quote}
Before the new one: say the Mandarin for new, then the Mandarin for fast.
\end{quote}

\subsection*{You'll want to know: jiù}
\emph{jiù} — \textquotedblleft{}old (of things); no longer new\textquotedblright{}.

\begin{quote}
\emph{wǒ de chē hěn jiù} — My car is old.
\end{quote}

Its characters come later; for now it is a word for the ear.

\begin{cousinweb}
\emph{jiù} is inherited from Middle Chinese \emph{gjuwH} (Guangyun, in Baxter's transcription), glossed 'former, of old'.
\end{cousinweb}

\begin{grammarlens}[title={What you've built}]
One more word for everyday answers.
\end{grammarlens}

\subsection*{Your turn}
\begin{itemize}
\raggedright
  \item \emph{Say it:} \emph{jiù}
  \item \emph{Say it:} \emph{jiù}, once more
  \item \emph{Say it:} \emph{wǒ de chē hěn jiù} — My car is old.
  \item \emph{Recall:} say the Mandarin for new, then the Mandarin for fast, then say \emph{jiù} again
\end{itemize}

\begin{tcolorbox}[breakable,colback=teal!4,colframe=teal!35!black,title={Before you move on}]
What does jiù mean? (\textquotedblleft{}Old (of things); no longer new\textquotedblright{}.) Say it once more.
\end{tcolorbox}
\section[yuǎn]{yuǎn — far}

\label{lesson:ZH-C62-yuan}



\begin{quote}
Before the new one: say the Mandarin for fast, then the Mandarin for no longer new.
\end{quote}

\subsection*{You'll want to know: yuǎn}
\emph{yuǎn} — \textquotedblleft{}far\textquotedblright{}.

\begin{quote}
\emph{Běijīng hěn yuǎn} — Beijing is a long way away.
\end{quote}

Its characters come later; for now it is a word for the ear.

\begin{cousinweb}
\emph{yuǎn} is inherited from Middle Chinese \emph{hjwonX} (Guangyun, in Baxter's transcription), glossed 'distant'.
\end{cousinweb}

\begin{grammarlens}[title={What you've built}]
One more word for everyday answers.
\end{grammarlens}

\subsection*{Your turn}
\begin{itemize}
\raggedright
  \item \emph{Say it:} \emph{yuǎn}
  \item \emph{Say it:} \emph{yuǎn}, once more
  \item \emph{Say it:} \emph{Běijīng hěn yuǎn} — Beijing is a long way away.
  \item \emph{Recall:} say the Mandarin for fast, then the Mandarin for no longer new, then say \emph{yuǎn} again
\end{itemize}

\begin{tcolorbox}[breakable,colback=teal!4,colframe=teal!35!black,title={Before you move on}]
What does yuǎn mean? (\textquotedblleft{}Far\textquotedblright{}.) Say it once more.
\end{tcolorbox}
\section[jìn]{jìn — near; close}

\label{lesson:ZH-C62-jin}



\begin{quote}
Before the new one: say the Mandarin for no longer new, then the Mandarin for far.
\end{quote}

\subsection*{You'll want to know: jìn}
\emph{jìn} — \textquotedblleft{}near; close\textquotedblright{}.

\begin{quote}
\emph{wǒ jiā hěn jìn} — My home is close by.
\end{quote}

Its characters come later; for now it is a word for the ear.

\begin{cousinweb}
\emph{jìn} is inherited from Middle Chinese \emph{gj+nX} (Guangyun, in Baxter's transcription), glossed 'pressing close' and 'almost'.
\end{cousinweb}

\begin{grammarlens}[title={What you've built}]
One more word for everyday answers.
\end{grammarlens}

\subsection*{Your turn}
\begin{itemize}
\raggedright
  \item \emph{Say it:} \emph{jìn}
  \item \emph{Say it:} \emph{jìn}, once more
  \item \emph{Say it:} \emph{wǒ jiā hěn jìn} — My home is close by.
  \item \emph{Recall:} say the Mandarin for no longer new, then the Mandarin for far, then say \emph{jìn} again
\end{itemize}

\begin{tcolorbox}[breakable,colback=teal!4,colframe=teal!35!black,title={Before you move on}]
What does jìn mean? (\textquotedblleft{}Near; close\textquotedblright{}.) Say it once more.
\end{tcolorbox}
\section[zhīdào]{zhīdào — to know a fact; to be aware}

\label{lesson:ZH-C62-zhidao}



\begin{quote}
Before the new one: say the Mandarin for far, then the Mandarin for near.
\end{quote}

\subsection*{You'll want to know: zhīdào}
\emph{zhīdào} — \textquotedblleft{}to know a fact; to be aware\textquotedblright{}.

\begin{quote}
\emph{wǒ zhīdào nǐ xiǎng shuō shénme} — I know what you want to say.
\end{quote}

Its characters come later; for now it is a word for the ear.

\begin{cousinweb}
\emph{zhīdào} is built from zhī 'to perceive' and dào 'the way', both inherited from Middle Chinese, \emph{trje} and \emph{dawX} (Guangyun, in Baxter's transcription). How the pair came to mean knowing a fact is not settled.
\end{cousinweb}

\begin{grammarlens}[title={What you've built}]
One more word for everyday answers.
\end{grammarlens}

\subsection*{Your turn}
\begin{itemize}
\raggedright
  \item \emph{Say it:} \emph{zhīdào}
  \item \emph{Say it:} \emph{zhīdào}, once more
  \item \emph{Say it:} \emph{wǒ zhīdào nǐ xiǎng shuō shénme} — I know what you want to say.
  \item \emph{Recall:} say the Mandarin for far, then the Mandarin for near, then say \emph{zhīdào} again
\end{itemize}

\begin{tcolorbox}[breakable,colback=teal!4,colframe=teal!35!black,title={Before you move on}]
What does zhīdào mean? (\textquotedblleft{}To know a fact; to be aware\textquotedblright{}.) Say it once more.
\end{tcolorbox}
